\section{Methods: Group 5 --- Signal and Shape Screens}
\label{sec:g5}

Group 5 screens for abnormal findings the geometric groups do not: vertebral-body
compression/deformity and intramedullary cord signal change (myelomalacia). Both are screens that flag
findings for physician review, never diagnoses.

\paragraph{5.2 Vertebral compression/deformity screen.} This reuses the Group~1 \hahp{} machinery
(canal-cut body isolation $+$ endplate-line heights) and flags \hahp{} below the cohort-calibrated cut
($\approx 0.68$, mean $-2$SD). Scope honesty is essential: validated against RSNA-2022 cervical fractures,
the geometric wedge metric has near-zero power on \emph{non-compression} fractures (odontoid/facet/arch),
so 5.2 is scoped explicitly as a vertebral-body \emph{compression} screen, not a general fracture
detector. Healthy false-positive rate fell from 17\% to 0\% after recalibration
(Section~\ref{sec:journey}).

\paragraph{5.1 Myelomalacia screen.} Hand-rolled intensity thresholds cannot separate lesion from normal
cord signal variation, so we adopt SCIseg~\cite{sciseg2024} (\texttt{sct\_deepseg lesion\_sci\_t2}) as the
engine; sensitivity rests on its publication. Our contribution is the \emph{specificity} check on healthy
cords: 10/11 healthy cords were completely clean (the one false positive a small component at the
cervicothoracic junction), $\approx$91\% healthy specificity in our hands. Lesions are mapped to the
cervical level they overlap and carried into the report.

\paragraph{5.3 / 5.4 scoped out.} Tumour/mass detection (5.3) is scoped out --- no public labelled
cervical-tumour MRI exists to validate it. Post-surgical scar (5.4) is deferred --- distinguishing scar
from recurrent disc requires gadolinium-enhanced sequences not in the T2-only input. Both are listed under
``not assessed'' in the output contract.

\paragraph{The 5$\rightarrow$6 contract.} Group 5 emits a per-level findings document (compression screen
$+$ myelomalacia, each with value, status, and citable provenance, plus an explicit not-assessed list and
no-diagnosis wording) that the interpretation layer consumes.

\paragraph{Status.} Group 5 is validated: 5.2 (FP 17\%$\rightarrow$0\%) and 5.1 ($\approx$91\% healthy
specificity), with the compression scope and the no-MRI-comparator caveats documented.
